% =========================================================
%  PLANTILLA DE MATRIZ DE RIESGOS
%  Actividad A03 - Administración de Riesgos
%  Seguridad Aplicada a Sistemas de Información
%  Compilar con: pdfLaTeX (Overleaf o TeX Live/MiKTeX)
% =========================================================
\documentclass[11pt,a4paper]{article}

% ----------------- Paquetes -----------------
\usepackage[utf8]{inputenc}
\usepackage[spanish]{babel}
\usepackage{geometry}
\geometry{left=2.2cm,right=2.2cm,top=2.5cm,bottom=2.5cm}
\usepackage[table]{xcolor}
\usepackage{tcolorbox}
\tcbuselibrary{skins,breakable}
\usepackage{enumitem}
\usepackage{tabularx}
\usepackage{booktabs}
\usepackage{titlesec}
\usepackage{fancyhdr}
\usepackage{hyperref}
\usepackage{graphicx}
\usepackage{multirow}
\usepackage{array}

% ----------------- Colores institucionales -----------------
\definecolor{azulmat}{RGB}{0,84,140}
\definecolor{griscaja}{RGB}{245,245,245}
\definecolor{azulsuave}{RGB}{232,240,250}
\definecolor{amarillosuave}{RGB}{255,255,224}
% Colores de la matriz de calor
\definecolor{nivelbajo}{RGB}{144,238,144}
\definecolor{nivelmedio}{RGB}{255,235,120}
\definecolor{nivelalto}{RGB}{255,170,90}
\definecolor{nivelcritico}{RGB}{255,90,90}

\hypersetup{
  colorlinks=true,
  linkcolor=azulmat,
  urlcolor=azulmat,
  citecolor=azulmat
}

% ----------------- Cajas de contenido -----------------
\newtcolorbox{instruccion}[1][Instrucciones]{%
  enhanced, breakable,
  colback=amarillosuave, colframe=orange!75!black,
  fonttitle=\bfseries, title=#1
}

\newtcolorbox{informacion}[1][Información]{%
  enhanced, breakable,
  colback=azulsuave, colframe=azulmat,
  fonttitle=\bfseries, title=#1
}

% Caja vacía para completar (texto libre)
\newtcolorbox{campolibre}[1][2cm]{%
  enhanced, breakable,
  colback=white, colframe=griscaja!80!black,
  height=#1, valign=top, boxrule=0.5pt
}

% ----------------- Encabezado y pie -----------------
\pagestyle{fancy}
\fancyhf{}
\fancyhead[L]{\small Plantilla de Matriz de Riesgos --- A03}
\fancyhead[R]{\small Seguridad Aplicada a Sistemas de Información}
\fancyfoot[C]{\thepage}
\renewcommand{\headrulewidth}{0.4pt}

% ----------------- Formato de títulos -----------------
\titleformat{\section}{\large\bfseries\color{azulmat}}{\thesection.}{0.6em}{}
\titleformat{\subsection}{\normalsize\bfseries\color{azulmat!80!black}}{\thesubsection.}{0.6em}{}

% ----------------- Portada -----------------
\title{%
  \rule{\textwidth}{1pt}\\[0.4cm]
  \textbf{\Large Plantilla de Matriz de Riesgos}\\[0.2cm]
  \large Actividad A03 --- Administración de Riesgos\\[0.1cm]
  \normalsize Análisis y Gestión de Riesgos en Sistemas de Información\\[0.2cm]
  \rule{\textwidth}{1pt}\\[0.4cm]
  \normalsize Seguridad Aplicada a Sistemas de Información\\
  Licenciatura en Sistemas de Información --- Facultad de Informática y Diseño\\
  Docente: Ing.\ Rodrigo Atilio Elgueta\\[0.2cm]
  \small Ciclo Académico Vigente
}
\author{}
\date{}

\begin{document}
\maketitle
\thispagestyle{fancy}
\tableofcontents
\newpage

% =========================================================
\begin{instruccion}[Cómo usar esta plantilla]
Esta plantilla es el instrumento formal para desarrollar la \textbf{Actividad A03 --- Administración de Riesgos}. Completá cada sección con la información de la empresa real, ficticia o caso del Trabajo Integrador que hayas elegido.

\textbf{Recomendaciones:}
\begin{itemize}[leftmargin=*]
  \item Completá los campos de la Sección~\ref{sec:datos} antes de iniciar el análisis.
  \item Asigná un \textbf{código (ID)} único a cada activo, amenaza y riesgo para poder referenciarlos en las tablas.
  \item Usá las escalas de las Secciones~\ref{sec:escalas} para valorar probabilidad e impacto de forma consistente.
  \item Para cada riesgo identificado, calculá el \textbf{valor de riesgo} (\(Probabilidad \times Impacto\)) y clasificá el nivel según la Sección~\ref{sec:niveles}.
  \item Proponé salvaguardas y calculá el \textbf{riesgo residual} (Sección~\ref{sec:tratamiento}).
  \item Subí el informe completado a la carpeta del Trabajo Integrador en Google Drive, otorgando permisos de comentario a \texttt{era.profe@gmail.com}.
\end{itemize}
\end{instruccion}

% =========================================================
\section{Datos Generales}\label{sec:datos}

\begin{center}
\renewcommand{\arraystretch}{1.6}
\begin{tabularx}{\textwidth}{@{} p{5cm} X @{}}
\toprule
\textbf{Nombre de la empresa / organización} & \\
\midrule
\textbf{Rubro / Industria} & \\
\midrule
\textbf{Tamaño (empleados / facturación)} & \\
\midrule
\textbf{Alumno(s) / Grupo} & \\
\midrule
\textbf{Legajo(s)} & \\
\midrule
\textbf{Fecha de elaboración} & \\
\midrule
\textbf{Responsable del análisis} & \\
\bottomrule
\end{tabularx}
\end{center}

% =========================================================
\section{Objetivo y Alcance del Análisis}

\begin{campolibre}[2.5cm]
\end{campolibre}

\textbf{Alcance:} (sistemas, procesos, sedes o áreas incluidas en el análisis)

\begin{campolibre}[2cm]
\end{campolibre}

% =========================================================
\section{Inventario y Clasificación de Activos}

Asigná un código único a cada activo (ej.: A01, A02\ldots). La clasificación de la información puede ser: \textit{Pública, Interna, Confidencial, Restringida}.

\begin{center}
\small
\renewcommand{\arraystretch}{1.6}
\begin{tabularx}{\textwidth}{@{} >{\centering\arraybackslash}p{1cm} X >{\raggedright\arraybackslash}p{2.2cm} >{\raggedright\arraybackslash}p{2.3cm} >{\raggedright\arraybackslash}p{2.3cm} >{\centering\arraybackslash}p{1.8cm} @{}}
\toprule
\textbf{ID} & \textbf{Descripción del Activo} & \textbf{Tipo} & \textbf{Responsable} & \textbf{Clasificación} & \textbf{Criticidad} \\
\midrule
 & & & & & \\
 & & & & & \\
 & & & & & \\
 & & & & & \\
 & & & & & \\
 & & & & & \\
 & & & & & \\
 & & & & & \\
 & & & & & \\
 & & & & & \\
\bottomrule
\end{tabularx}
\end{center}

\textit{Tipos de activo: Información / Software / Hardware / Red / Humano / Imagen. Criticidad: Baja / Media / Alta.}

% =========================================================
\section{Identificación de Amenazas y Vulnerabilidades}

Asigná un código único a cada amenaza (ej.: T01, T02\ldots). El tipo puede ser: \textit{Natural, Accidental, Intencional}.

\begin{center}
\small
\renewcommand{\arraystretch}{1.6}
\begin{tabularx}{\textwidth}{@{} >{\centering\arraybackslash}p{1cm} >{\centering\arraybackslash}p{1.8cm} X >{\raggedright\arraybackslash}p{2.2cm} X @{}}
\toprule
\textbf{ID} & \textbf{Activo (ID)} & \textbf{Descripción de la Amenaza} & \textbf{Tipo} & \textbf{Vulnerabilidad Asociada} \\
\midrule
 & & & & \\
 & & & & \\
 & & & & \\
 & & & & \\
 & & & & \\
 & & & & \\
 & & & & \\
 & & & & \\
\bottomrule
\end{tabularx}
\end{center}

% =========================================================
\section{Escalas de Valoración}\label{sec:escalas}

\subsection{Escala de Probabilidad}

\begin{center}
\renewcommand{\arraystretch}{1.4}
\begin{tabularx}{\textwidth}{@{} c p{2.8cm} X @{}}
\toprule
\textbf{Valor} & \textbf{Nivel} & \textbf{Descripción} \\
\midrule
1 & Raro & El evento solo ocurriría en circunstancias excepcionales. \\
\midrule
2 & Improbable & Podría ocurrir, pero no se espera que suceda. \\
\midrule
3 & Posible & Existe una posibilidad real de que ocurra. \\
\midrule
4 & Probable & Es muy probable que ocurra en algún momento. \\
\midrule
5 & Casi seguro & Se espera que ocurra frecuentemente / ha ocurrido recientemente. \\
\bottomrule
\end{tabularx}
\end{center}

\subsection{Escala de Impacto}

\begin{center}
\renewcommand{\arraystretch}{1.4}
\begin{tabularx}{\textwidth}{@{} c p{2.8cm} X @{}}
\toprule
\textbf{Valor} & \textbf{Nivel} & \textbf{Descripción} \\
\midrule
1 & Insignificante & Impacto mínimo, sin consecuencias operativas ni económicas relevantes. \\
\midrule
2 & Menor & Alteración leve de la operación, costo bajo. \\
\midrule
3 & Moderado & Impacto operativo y económico apreciable, recuperable. \\
\midrule
4 & Mayor & Impacto significativo, pérdida de operación o datos, costo alto. \\
\midrule
5 & Catastrófico & Paralización total, pérdida crítica de datos, daño reputacional o legal severo. \\
\bottomrule
\end{tabularx}
\end{center}

% =========================================================
\section{Matriz de Calor (Probabilidad × Impacto)}

La matriz 5×5 relaciona la \textbf{probabilidad} (eje vertical) con el \textbf{impacto} (eje horizontal). Cada celda muestra el valor de riesgo resultante (\(P \times I\)) y su color indica el nivel.

\begin{center}
\renewcommand{\arraystretch}{1.9}
\setlength{\tabcolsep}{8pt}
\begin{tabular}{@{} c c | c c c c c @{}}
\multicolumn{2}{c}{} & \multicolumn{5}{c}{\textbf{IMPACTO} $\longrightarrow$} \\
\multicolumn{2}{c}{} & \textbf{1} & \textbf{2} & \textbf{3} & \textbf{4} & \textbf{5} \\
\cline{3-7}
\multirow{5}{*}{\rotatebox[origin=c]{90}{\textbf{PROBABILIDAD} $\longrightarrow$}}
 & \textbf{5} & \cellcolor{nivelmedio}5 & \cellcolor{nivelalto}10 & \cellcolor{nivelalto}15 & \cellcolor{nivelcritico}20 & \cellcolor{nivelcritico}25 \\
 & \textbf{4} & \cellcolor{nivelbajo}4 & \cellcolor{nivelmedio}8 & \cellcolor{nivelalto}12 & \cellcolor{nivelcritico}16 & \cellcolor{nivelcritico}20 \\
 & \textbf{3} & \cellcolor{nivelbajo}3 & \cellcolor{nivelmedio}6 & \cellcolor{nivelmedio}9 & \cellcolor{nivelalto}12 & \cellcolor{nivelalto}15 \\
 & \textbf{2} & \cellcolor{nivelbajo}2 & \cellcolor{nivelbajo}4 & \cellcolor{nivelmedio}6 & \cellcolor{nivelmedio}8 & \cellcolor{nivelalto}10 \\
 & \textbf{1} & \cellcolor{nivelbajo}1 & \cellcolor{nivelbajo}2 & \cellcolor{nivelbajo}3 & \cellcolor{nivelbajo}4 & \cellcolor{nivelmedio}5 \\
\cline{3-7}
\end{tabular}
\end{center}

\subsection{Niveles de Riesgo}\label{sec:niveles}

\begin{center}
\renewcommand{\arraystretch}{1.4}
\begin{tabularx}{\textwidth}{@{} p{3cm} p{2.5cm} p{3.5cm} X @{}}
\toprule
\textbf{Nivel} & \textbf{Rango (P×I)} & \textbf{Color} & \textbf{Acción requerida} \\
\midrule
\cellcolor{nivelbajo} \textbf{Bajo} & 1 -- 4 & Verde & Monitorear. Puede aceptarse. \\
\cellcolor{nivelmedio} \textbf{Medio} & 5 -- 9 & Amarillo & Plan de acción a mediano plazo. \\
\cellcolor{nivelalto} \textbf{Alto} & 10 -- 15 & Naranja & Tratamiento prioritario a corto plazo. \\
\cellcolor{nivelcritico} \textbf{Crítico} & 16 -- 25 & Rojo & Acción inmediata. Escalar a dirección. \\
\bottomrule
\end{tabularx}
\end{center}

% =========================================================
\section{Evaluación de Riesgos}\label{sec:evaluacion}

Asigná un código único a cada riesgo (ej.: R01, R02\ldots). Calculá el valor como \(Probabilidad \times Impacto\) y clasificá el nivel según la tabla anterior.

\begin{center}
\small
\renewcommand{\arraystretch}{1.6}
\begin{tabularx}{\textwidth}{@{} >{\centering\arraybackslash}p{0.8cm} >{\centering\arraybackslash}p{1.2cm} >{\centering\arraybackslash}p{1.6cm} >{\centering\arraybackslash}p{1.1cm} >{\centering\arraybackslash}p{1.3cm} >{\centering\arraybackslash}p{1.1cm} >{\centering\arraybackslash}p{1.4cm} X @{}}
\toprule
\textbf{ID} & \textbf{Activo} & \textbf{Amenaza} & \textbf{Prob.} & \textbf{Impacto} & \textbf{Valor} & \textbf{Nivel} & \textbf{Justificación breve} \\
\midrule
 & & & & & & & \\
 & & & & & & & \\
 & & & & & & & \\
 & & & & & & & \\
 & & & & & & & \\
 & & & & & & & \\
 & & & & & & & \\
 & & & & & & & \\
\bottomrule
\end{tabularx}
\end{center}

% =========================================================
\section{Tratamiento de Riesgos y Riesgo Residual}\label{sec:tratamiento}

Para cada riesgo, definí la estrategia de tratamiento:
\begin{itemize}[leftmargin=*]
  \item \textbf{Mitigar:} implementar salvaguardas para reducir probabilidad o impacto.
  \item \textbf{Transferir:} pasar el riesgo a un tercero (seguros, outsourcing).
  \item \textbf{Aceptar:} asumir el riesgo conscientemente.
  \item \textbf{Evitar:} discontinuar la actividad que lo genera.
\end{itemize}

\begin{center}
\small
\renewcommand{\arraystretch}{1.6}
\begin{tabularx}{\textwidth}{@{} >{\centering\arraybackslash}p{1.3cm} >{\raggedright\arraybackslash}p{1.7cm} X >{\raggedright\arraybackslash}p{1.7cm} >{\centering\arraybackslash}p{1cm} >{\centering\arraybackslash}p{1cm} >{\centering\arraybackslash}p{1cm} >{\centering\arraybackslash}p{1.3cm} @{}}
\toprule
\textbf{Riesgo} & \textbf{Estrategia} & \textbf{Salvaguardas propuestas} & \textbf{Tipo salv.} & \textbf{Prob. resid.} & \textbf{Imp. resid.} & \textbf{Val. resid.} & \textbf{Nivel resid.} \\
\midrule
 & & & & & & & \\
 & & & & & & & \\
 & & & & & & & \\
 & & & & & & & \\
 & & & & & & & \\
 & & & & & & & \\
\bottomrule
\end{tabularx}
\end{center}

\textit{Tipo de salvaguarda: Técnica / Física / Administrativa. El nivel residual se clasifica con la misma tabla de la Sección~\ref{sec:niveles}.}

% =========================================================
\section{Resumen de Resultados}

Distribución de los riesgos identificados por nivel (completar con la cantidad de riesgos en cada categoría):

\begin{center}
\renewcommand{\arraystretch}{1.4}
\begin{tabularx}{0.85\textwidth}{@{} p{4cm} p{3.5cm} X @{}}
\toprule
\textbf{Nivel} & \textbf{Cantidad de riesgos} & \textbf{\% del total} \\
\midrule
\cellcolor{nivelcritico} Crítico & & \\
\cellcolor{nivelalto} Alto & & \\
\cellcolor{nivelmedio} Medio & & \\
\cellcolor{nivelbajo} Bajo & & \\
\midrule
\textbf{Total} & & 100\% \\
\bottomrule
\end{tabularx}
\end{center}

\section{Conclusiones y Recomendaciones}

\begin{campolibre}[4cm]
\end{campolibre}

\section{Declaración de Buenas Prácticas}

\begin{informacion}[Compromiso profesional]
Como estudiante de Seguridad Aplicada a Sistemas de Información, declaro que el presente análisis de riesgos fue elaborado aplicando las buenas prácticas de la disciplina, con criterio ético, fundamentos técnicos y respetando el marco normativo vigente.
\end{informacion}

\vspace{1.2cm}

\begin{center}
\begin{tabularx}{\textwidth}{@{} X X @{}}
\rule{6cm}{0.4pt} & \rule{6cm}{0.4pt} \\
Firma del alumno/a & Fecha \\
\end{tabularx}
\end{center}

% =========================================================
\section{Referencias Bibliográficas}

\begin{itemize}[leftmargin=*]
  \item Chicano Tejada, E. (2023). \textit{Auditoría de seguridad informática. IFCT0109} (2.ª ed.). IC Editorial --- eLibro.
  \item ISO/IEC (2005). \textit{Information security management 27001}.
  \item ISO/IEC (2018). \textit{ISO/IEC 27005: Information security risk management}.
  \item NIST (2012). \textit{Guide for Conducting Risk Assessments (SP 800-30 Rev. 1)}.
  \item OWASP Risk Rating Methodology: \url{https://owasp.org/www-community/OWASP_Risk_Rating_Methodology}
\end{itemize}

\end{document}